\chapter{The original evaluation, reproduced from source}
\label{ch:original}

\begin{saybox}[title=What to say at the start of this section]
State that the original protocol was not merely read but re-executed, with
the production seed, so the numbers discussed are the actual production
numbers rather than a description of them. Then present the protocol's
design intent honestly: it was written to prevent a specific, real form of
leakage, and it succeeds at that.
\end{saybox}

\section{The splitting procedure as implemented}

The production split is performed by \texttt{make\_context\_splits()} in
\evd{model\_service.py}, called from the training driver with the global
seed. Its logic, read from source:

\begin{enumerate}[leftmargin=1.5em, itemsep=2pt]
\item Compute, for each \texttt{context\_id}, whether that context contains
  any row with \texttt{data\_origin == "real\_paper\_derived"}.
\item Partition the context identifiers into two strata by that boolean.
\item Within each stratum, shuffle the context identifiers with
  \texttt{numpy.random.default\_rng(42)} and allocate them $70/15/15$ to
  train, validation and test according to
  \texttt{SPLIT\_FRACTIONS = \{train: 0.70, validation: 0.15, test: 0.15\}}.
\item Return the three row subsets selected by context membership.
\end{enumerate}

The design intent is explicit in the function's own docstring: it is a
``leakage-safe split'' that groups by context ``so a context is never split
across train/val/test,'' stratified so that a small real-data subset would
be represented in every partition rather than left to chance.

\begin{findingbox}[title=The protocol does what it claims --- and that is the point]
Re-executing the production splitter with the production seed and checking
context membership directly confirms \textbf{zero context overlap between
train and test in all six specialists}. The protocol is not broken. It
correctly enforces the constraint it was designed to enforce. The
investigation's finding is not that the code is wrong, but that the
constraint it enforces is not the constraint this dataset requires.
\end{findingbox}

\subsection{An observation about the stratification}

The stratification variable is \texttt{data\_origin ==
"real\_paper\_derived"}. In the V7 corpus, every row has
\texttt{data\_origin == "synthetic\_literature\_constrained\_v7"}. The
predicate is therefore uniformly false, one stratum is empty, and the
stratification is a no-op on this data.

This is not an error --- it is a piece of machinery written when real rows
were expected to be present, and it would function if they were. It is
worth mentioning to the panel because it is independent corroboration of
\Cref{sec:limit-evidence}: the pipeline was built anticipating real
observations that the corpus does not currently contain.

\section{The original numerical results}

\Cref{tab:tab_production_v7} gives the complete production results for the
V7 dataset: all six specialists $\times$ four model families, with the
training, validation and test $\rsq$ and the test RMSE, recovered from the
stored \texttt{metrics.json} and \texttt{training\_manifest.json} artefacts
rather than from any summary document. \Cref{tab:tab_production_v6} gives
the corresponding V6 results.

\tbl{tab_production_v7}{Original production results, V7 dataset,
context-holdout protocol. Recovered from
\texttt{artifacts\_v7/\{category\}/\{task\}/metrics.json}.}

\fig{fig01_production_v7_test_r2}{0.95}{%
  Original context-holdout test $\rsq$ for the twenty-four production fits.
  The three general-shelf-life specialists produce the scores that prompted
  this investigation; the safety-endpoint specialists behave very
  differently and are discussed in \Cref{sec:safety-anomaly}.}

\subsection{Reading the table honestly}

Three distinct patterns are present, and conflating them would be a
mistake.

\paragraph{The general-shelf-life specialists score very high.}
Soft/general reaches test $\rsq = 0.971$ (LightGBM), hard/general $0.932$,
semi-hard/general $0.927$. These are the numbers that prompted the concern.

\paragraph{The safety-endpoint specialists do not.}
\label{sec:safety-anomaly}
On V7, soft/safety reaches only $0.594$, semi-hard/safety $0.342$, and
hard/safety is \emph{negative} at $-0.097$ for Random Forest and
$-0.528$ for EBM. A negative $\rsq$ means the model predicts worse than
the test-set mean. Any explanation of the high scores must also be
consistent with these low ones, which rules out ``the models memorise
everything'' as a general account.

\paragraph{The V6\,$\rightarrow$\,V7 transition moved the two tasks in
opposite directions.} Comparing \Cref{tab:tab_production_v6} with
\Cref{tab:tab_production_v7}: soft/general improved from $0.936$ to
$0.971$, while soft/safety \emph{fell} from $0.917$ to $0.573$ and
hard/safety fell from $0.786$ to $-0.024$ (LightGBM in each case).

\fig{fig02_v6_vs_v7_lightgbm}{0.95}{%
  The same model and the same protocol on two dataset versions. The
  divergence between quality and safety endpoints under V7 is a property
  of the data, not of the learner.}

\tbl{tab_production_v6}{Original production results, V6 dataset, same
protocol. Retained for comparison; not overwritten by the V7 run.}

\subsection{Optimism gap}

\Cref{fig:fig03_train_test_gap_v7} plots the difference between training
and test $\rsq$ for every V7 production fit. Random Forest and XGBoost
train to $\rsq > 0.99$ almost everywhere, which is expected for
unpruned bagged trees and deep boosted trees --- a training $\rsq$ near
unity is a statement about capacity, not about generalization. The
informative quantity is the gap, and the gap is small precisely where the
test score is high.

\fig{fig03_train_test_gap_v7}{0.95}{%
  Train-minus-test $\rsq$ for the V7 production fits. The large gaps sit
  on the safety-endpoint specialists, where the test scores are low; the
  quality specialists show small gaps and high test scores.}

\section{Four things a high $\rsq$ could mean}

The single most important conceptual point of the meeting is that these
four explanations are distinct, and the original result is consistent with
only some of them. They must not be used interchangeably.

\begin{description}[leftmargin=1.6em, style=nextline, font=\bfseries\color{dgreen}]
\item[1. Classical overfitting] The model has memorised individual
  training rows, including their noise, and performs well on training data
  and poorly on any held-out data. \emph{This is refuted by the original
  results themselves}: held-out test scores are high, not low. Whatever is
  happening, it is not the textbook train-high/test-low failure.

\item[2. Data leakage] A feature carries information about the target that
  would not be available at prediction time --- an identifier, a downstream
  measurement, an encoded target. \emph{Directly checked}: the exclusion
  list removes identifiers and provenance, including
  \texttt{source\_rule\_id} itself, and \Cref{sec:leakage} audits the
  remaining features individually. No target-derived feature was found in
  the model inputs.

\item[3. Optimistic evaluation from shared generative structure] The test
  set is statistically independent at the level the protocol controls
  (contexts) but not at the level that governs the mapping being learned
  (generation rules). The model is scored on new samples of relationships
  it was trained on. \emph{This is what the evidence supports}, and
  \Cref{ch:discovery} quantifies it.

\item[4. Genuine generalization to independent observations] The model
  would predict well on new experimental data. \emph{Not established by
  any internal metric}, and \Cref{ch:external} shows the external evidence
  points the other way.
\end{description}

\begin{findingbox}[title=The precise claim]
The original $\rsq \approx 0.92$--$0.97$ figures are best described as
\textbf{valid measurements of the wrong quantity}. They accurately measure
how well the model interpolates within known generation mechanisms. They
were reported, and read, as if they measured generalization to new
formulations. The fault is in the inference drawn from the protocol, not
in the arithmetic of the metric.
\end{findingbox}
